\chapter{Backend Architecture}
\label{ch:backend-architecture}

Every backend chapter that follows implements the same contract. This chapter covers that
contract once --- the interfaces, the shared parameter struct, and the build-time selection
mechanism --- so the nine backend-specific chapters ahead can spend their time on what
actually differs, not on repeating the shape they all share.

\section{One header, not several}

Despite \cnaclass{IGraphicsBackend} being described throughout this book as ``the backend
contract layer,'' it is worth being precise: the entire contract lives in a single header,
\texttt{include/CNA/Internal/Backends/Common/IGraphicsBackend.hpp} --- roughly a thousand
lines, not a family of separate \texttt{ISpriteBatchBackend.hpp}/\texttt{ITextureBackend.hpp}
files. Eleven interface classes live inside it, all in
\texttt{namespace CNA::Internal::Backends}:

\begin{center}
\begin{tabular}{ll}
\toprule
\textbf{Interface} & \textbf{Owns} \\
\midrule
\texttt{IVertexBufferBackend} & Vertex data upload, streaming hints, declaration \\
\texttt{IIndexBufferBackend} & 16/32-bit index data upload \\
\texttt{IOcclusionQueryBackend} & Begin/End/IsComplete/PixelCount \\
\texttt{ITextureCubeBackend} & Per-face texture data \\
\texttt{ITexture3DBackend} & Per-slice volume texture data \\
\texttt{ITextureBackend} & 2D texture data, GL interop, CPU shadow \\
\texttt{IRenderTargetBackend} & Extends \texttt{ITextureBackend}: bind/unbind as target \\
\texttt{IRenderTargetCubeBackend} & Extends \texttt{ITextureCubeBackend}: per-face target binding \\
\texttt{IEffectBackend} & Shader program compile/bind, uniform/texture setters \\
\texttt{ISpriteBatchBackend} & Begin/End, transform, three \texttt{Draw} overloads \\
\texttt{IGraphicsBackend} & Everything else --- see below \\
\bottomrule
\end{tabular}
\end{center}

\cnaclass{IGraphicsBackend} itself is the large one: device lifecycle
(\texttt{Clear}/\texttt{Present}/viewport and presentation-mode management), native handle
access (\texttt{GetWindowInternal()}/\texttt{GetRendererInternal()} --- both pure virtual,
and both carry a comment flagging the \texttt{SDL\_Window*}/\texttt{SDL\_Renderer*} return
types as ``should be abstracted later''), resource factories for every GPU object type in
Part~III, render-target binding, per-state-object apply methods, the full 3D draw dispatch
(\texttt{DrawColoredPrimitives}, the effect-aware \texttt{DrawPrimitivesEx} family, instanced
draws), and a capability-query mechanism (\S\ref{sec:graphicscapability} below). A single,
project-wide design choice worth internalizing before reading any backend chapter: nearly
every optional method on this interface \textbf{defaults to a safe no-op or a null return}
rather than requiring every backend to implement everything. \texttt{CreateEffectBackend()}
defaults to \texttt{nullptr} (``backends that do not support programmable shaders'');
\texttt{CreateOcclusionQuery()} defaults to \texttt{nullptr}; every state-application method
defaults to doing nothing at all. This is a deliberate ``graceful capability degradation''
design, and every backend chapter's own ``degrades/gap/blocked'' distinctions in Part~IV rest on it: a feature
that silently degrades to a no-op reads differently in an audit than one that throws, even
though both mean ``not implemented here.''

\section{GpuDrawParams: already covered, worth remembering here}

Chapter~\ref{ch:stock-effects} already introduced \cnaclass{GpuDrawParams} in full --- the
single struct \texttt{Effect::FillGpuDrawParams()} populates and every backend's
\texttt{DrawPrimitivesEx}/\texttt{DrawIndexedPrimitivesEx} consumes. It is worth restating
here only because it is the concrete mechanism by which ``a stock effect draws identically
across eleven backends'' is actually achieved: the effect layer never talks to a backend
directly about lights, fog, or skinning palettes --- it fills one struct, and each backend's
own draw dispatch decides how to turn that struct into real GPU state.

\section{GraphicsCapability: ask, don't catch}
\label{sec:graphicscapability}

Rather than requiring game code to attempt an operation and catch an exception to discover
whether a backend supports it, \cnaclass{GraphicsCapability} is a small enum ---
\texttt{ThreeD}, \texttt{DepthStencilBuffer}, \texttt{MultiSampleAntiAliasing},
\texttt{MultipleRenderTargets}, \texttt{AnisotropicFiltering}, \texttt{WireFrame},
\texttt{OcclusionQuery}, \texttt{CustomEffects} --- queryable via
\texttt{GraphicsDevice::SupportsCapability(GraphicsCapability)}. The default implementation
returns \texttt{true} for everything; a 2D-only backend like \texttt{SDL\_RENDERER}
(Chapter~\ref{ch:sdl-renderer}) overrides this to honestly report what it cannot do, rather
than accepting a 3D draw call and throwing only once the call actually arrives.

\section{Build-time selection}

\texttt{cmake/BackendSelection.cmake} is where the backend list from
Chapter~\ref{ch:building-cna} is actually enumerated: the \texttt{CNA\_GRAPHICS\_BACKEND}
cache variable's \texttt{STRINGS} property lists all fourteen recognized values
(\texttt{SDL\_RENDERER}, \texttt{EASYGL}, \texttt{BGFX}, \texttt{VULKAN}, \texttt{WEBGPU},
\texttt{HEADLESS}, \texttt{SOFTWARE}, \texttt{D3D11}, \texttt{D3D12}, \texttt{CANVAS},
\texttt{ASCII}, \texttt{DX3}, \texttt{D3D9}, and an additional experimental
\texttt{SDL\_GPU} not otherwise covered in this book), and thirteen corresponding boolean
options exist alongside it --- if more than one is somehow set \texttt{ON} at once,
configuration fails outright with ``Exactly one backend option must be ON.'' The default is
itself platform-conditional: \texttt{EASYGL} on Linux or Emscripten, \texttt{SDL\_RENDERER}
everywhere else.

Several backends carry hard, configure-time platform gates rather than soft warnings:
\texttt{D3D9}/\texttt{D3D11}/\texttt{D3D12} fail configuration outright unless
\texttt{CMAKE\_SYSTEM\_NAME} is \texttt{Windows} (native or MinGW-w64 cross-compiled) ---
the comment in the CMake script is blunt about why: ``d3d11.h/d3d12.h/dxgi.h do not exist
elsewhere.'' \texttt{CANVAS} fails configuration outright unless building under Emscripten,
since HTML \texttt{<canvas>} is a browser DOM API with no meaning anywhere else.
\texttt{EASYGL} and \texttt{BGFX} both emit a soft warning (not a hard failure) if configured
outside Linux. \texttt{EASYGL} and \texttt{DX3} both require their respective sibling
checkouts (\texttt{../easy-gl}, \texttt{../free-direct}) to exist at all, hard-erroring
otherwise; \texttt{BGFX} instead fetches \texttt{bkaradzic/bgfx.cmake} directly via CMake's
\texttt{FetchContent}.

\texttt{cmake/BackendLibraries.cmake} wires the resulting per-backend link libraries and
contains two architectural details worth knowing before Chapter~\ref{ch:canvas-ascii-backends}:
the \texttt{ASCII} backend is built by reusing \texttt{SDL\_RENDERER}'s own real \texttt{.cpp}
sources as a separate static library, compiled \emph{without} \texttt{CNA\_BACKEND\_SDL\_RENDERER}
defined so its own \texttt{CreateGraphicsBackend()} factory is compiled out and cannot clash
with ASCII's own factory --- confirming in build-system terms what
Chapter~\ref{ch:canvas-ascii-backends} states in prose: ASCII is composition around
SDL\_Renderer, not a reimplementation. And \texttt{D3D9}'s effect-compilation code
(\texttt{D3D9EffectBackend}/\texttt{D3D9ConstantTable}) is deliberately isolated into its own
static library specifically so that only that isolated library links against
\texttt{d3dcompiler} --- keeping the stock D3D9 rendering pipeline itself free of that
dependency, unlike \texttt{D3D11}/\texttt{D3D12}, which link \texttt{d3dcompiler}
unconditionally into their whole backend target.

\section{One factory function per backend}

Whichever backend is selected at configure time compiles exactly one implementation of
\texttt{std::unique\_ptr<IGraphicsBackend> CreateGraphicsBackend(const
GraphicsBackendCreateArgs\&)}, gated behind a \texttt{\#ifdef CNA\_BACKEND\_*} matching the
compile definition \texttt{BackendSelection.cmake} set. \cnaclass{GraphicsBackendCreateArgs}
itself carries the window handle, virtual resolution, presentation mode, context-recovery
flag, multisample count, swap interval, and --- as \texttt{NOXNA} additions introduced
alongside the \texttt{D3D9} backend --- back-buffer/depth-stencil format, fullscreen flag,
and graphics profile, each documented as a field ``every existing backend except D3D9 may
ignore'': a direct, structural acknowledgment that \texttt{D3D9} is the one backend that
actually consults these values the way a real device negotiation would.

With the shared contract established, the next nine chapters cover what each backend
actually does with it.
